\documentclass[10pt,a4paper]{amsart}
\usepackage[margin=19mm]{geometry}
\usepackage{amsmath,amssymb,mathtools}
\usepackage[scr=rsfs,cal=euler]{mathalfa}
\usepackage{xfrac}
\usepackage[unicode,hidelinks,bookmarks=false]{hyperref}
\newcommand{\ie}{i.e.\ }
\newcommand{\cf}{cf.\ }
\newcommand{\resp}{respectively\ }
\newcommand{\Subsection}{\subsection}
\providecommand{\nobreakdash}{\nobreak\hbox{-}}
\setlength{\emergencystretch}{2em}
\multlinegap0pt
\usepackage{amssymb}
\usepackage{bm}
\usepackage{enumerate}
\usepackage{tikz}
\usepackage{float}
\usepackage{caption}
\newtheorem{proposition}{Proposition}
\newtheorem{lemma}{Lemma}
\newtheorem{proofx}{Proofx}
\newcommand{\R}{\mathbb{R}}
\newcommand{\eps}{\varepsilon}
\newcommand{\opdiv}{\operatorname{div}}
\newcommand{\overbar}[1]{\mkern 1.5mu\overline{\mkern-1.5mu#1\mkern-1.5mu}\mkern 1.5mu}
\title{An $L^2$-quantitative global approximation \\ for the Stokes initial-boundary value problem}
\author{Mitsuo Higaki}
\date{Source: arXiv:2511.16079v2 (CC BY 4.0)}
\begin{document}
\maketitle

\noindent\emph{Source and licence.} An $L^2$-quantitative global approximation \\ for the Stokes initial-boundary value problem. Version arXiv:2511.16079v2 (CC BY 4.0), distributed by arXiv under the Creative Commons Attribution 4.0 International licence, http://creativecommons.org/licenses/by/4.0/, as recorded in the arXiv licence metadata for this exact version and shown on the abstract page https://arxiv.org/abs/2511.16079v2. Full licence notice: \texttt{licenses/arxiv-2511.16079-CC-BY-4.0.txt}.\par\smallskip

\noindent\textit{Editorial scope.} Counted unit: the article's proposition prop.global.approx (source0.tex lines 450--479). The article's complete original proof of it is delivered with it (source0.tex lines 1981--1983, one \texttt{proof} environment of 3 lines, which the article places away from the statement). No mathematics is written, repaired or re-derived: the statement and the proof are the article's own lines, in the article's own notation, and no retained line is substituted or re-flowed.  Retained uncounted as the source-local context the proof needs: lines 528--542; lines 667--693; lines 1023--1040; lines 1337--1350; lines 1379--1415; lines 1473--1480; lines 1490--1502; lines 1898--1910.  Nothing was omitted from any retained span, so the package carries no omission record.  No retained line contains a citation, so the wrapper carries no bibliography and no bibliography entry.  No retained line contains a figure environment, an image-inclusion command or drawing code, so no image is delivered and the package declares no asset dependency.  Editorial formatting only, in addition: an AMS \texttt{amsart} preamble that restates the article's own declarations and macros used by the retained lines, namely the environments \texttt{proposition}, \texttt{lemma}, \texttt{proofx} on separate counters, together with the standard packages those lines need.  The retained environments print in this document as Proposition 1, Lemma 1, Proofx 1 in order of appearance, whereas the article prints the same items with the numbering of its own full document.  The complete native source of the article as distributed in the arXiv e-print for this exact version is bundled unchanged as \texttt{sources/189492/source0.tex}.
\begin{lemma}[logarithmic stability]\label{lem.est.stab}
Let $(v, q)$ be given as in Proposition \ref{prop.global.approx}. Let $D'$ be a subdomain of $D$ satisfying $\overbar{D'} \subset D$. Suppose that, for $0<\eta<\mathcal{E}$, 
\[
    \|v\|_{H^1(D)} \le \mathcal{E}, 
    \qquad
    \|v\|_{L^2(D')} \le \eta. 
\]
Then, there exist $C>1$ and $\mu>0$ independent of $\lambda, \eta, \mathcal{E}$ such that 
\begin{equation}\label{est.stab}
    \|v\|_{L^2(D)}
    \le 
    C \mathcal{E} 
    \bigg(\log \frac{\mathcal{E}}{\eta} \bigg)^{-\mu}. 
\end{equation}
\end{lemma}

\begin{proposition}\label{prop.global.approx.step1}
Let $(v, q)$ be given as in Proposition \ref{prop.global.approx}. Then, for $\eps > 0$, there exists $F\in L^2(Y)^3$ such that 
\[
    v_1 
    := 
    \vec{E}_{\lambda} \ast {\mathbb P}_{\R^3} (\mathtt{e}_Y F) 
\]
approximates $v$ in $D$ as 
\[
    \|v - v_1\|_{L^2(D)}
    \le
    \eps 
    \Big( \|v\|_{H^{1}(D)} + \|q\|_{H^{1}(D)} \Big). 
\]
Moreover, $F$ is quantitatively estimated as 
\begin{equation}\label{est.F}
    \|F\|_{L^2(Y)}
    \le
    \exp
    \bigg(
    C
    \Big(\frac{\langle \lambda \rangle}{\eps} \Big)^{4/\mu}
    \bigg)
    \|v\|_{L^2(D)}, 
\end{equation}
where $C$ is independent of $\eps, \lambda$ while $\mu$ is introduced in Lemma \ref{lem.est.stab}. 
\end{proposition}

\begin{proposition}\label{prop.global.approx.step2}
Let $v_1$ be given as in Proposition \ref{prop.global.approx.step1}. Then, for $\eps>0$, there exists a smooth global approximation $u$ that solves 
\[
    \left\{
    \begin{array}{ll}
    \lambda u - \Delta u + \nabla p
    = 0&\mbox{in}\ \R^3, \\
    \opdiv u
    = 0&\mbox{in}\ \R^3, 
    \end{array}\right.
\]
for some smooth $p$, and approximates $v_1$ in $D$ as 
\[
    \|v_1 - u\|_{L^2(D)} 
    \le 
    \eps \|v\|_{L^2(D)}. 
\]
\end{proposition}

\begin{equation}\label{def.l0}
\begin{split}
    l_0
    = 
    \bigg(
    \frac{\eps^2}{C+1}
    \bigg)^{-1/3}
    \exp
    \bigg(
    \frac23
    \Big(\frac{\langle \lambda \rangle}{\eps} \Big)^{4/\mu}
    \bigg), 
\end{split}
\end{equation}

\begin{proposition}\label{prop.global.approx.step3}
Let $u$ be given as in Proposition \ref{prop.global.approx.step2}. Then, $u$ can be decomposed into 
\[
    u = u_1 + u_2, 
\]
where $u_1$ satisfies the Stokes resolvent problem, for some smooth $p_1$, 
\[
    \left\{
    \begin{array}{ll}
    \lambda u_1 - \Delta u_1 + \nabla p_1
    = 0&\mbox{in}\ \R^3, \\
    \opdiv u_1
    = 0&\mbox{in}\ \R^3 
    \end{array}\right.
\]
and $u_2$ the heat resolvent problem 
\[
    \left\{
    \begin{array}{ll}
    \lambda u_2 - \Delta u_2 
    = 0&\mbox{in}\ \R^3, \\
    \opdiv u_2
    = 0&\mbox{in}\ \R^3.
    \end{array}\right.
\]
Moreover, the following qualitative estimates hold: 
\begin{equation}\label{est.u1u2.qual}
    |u_1(x,\lambda)| 
    \lesssim
    \langle x \rangle^{l_0}, 
    \qquad
    |u_2(x,\lambda)| 
    \lesssim
    e^{(\Re\sqrt{\lambda}) |x|}, 
\end{equation}
where the implicit constants depend on $\eps, \lambda$ while $l_0$ is defined in \eqref{def.l0}.
\end{proposition}

\begin{equation}\label{def.N}
    N_{\eps,\lambda}
    = 
    C
    \Big(\frac{\langle \lambda \rangle}{\eps} \Big)^{4/\mu},
    \quad
    C>0,
\end{equation}

\begin{proposition}\label{prop.global.approx.step4}
Let $u_1$ be given as in Proposition \ref{prop.global.approx.step3}. Let $N_{\eps, \lambda}$ be defined in \eqref{def.N}. Then, $u_1$ is quantitatively estimated as 
\begin{equation}\label{est.u1}
\begin{split}
    |u_1(x,\lambda)|
    &\le
    \exp\big(\exp(N_{\eps,\lambda})\big)
    \|v\|_{L^2(D)} 
    \langle x \rangle^{\exp (N_{\eps,\lambda})}, 
\end{split}
\end{equation}
for sufficiently large $C$ independent of $\eps, \lambda$. 
\end{proposition}

\begin{proposition}\label{prop.global.approx.step5}
Let $u_2$ be given as in Proposition \ref{prop.global.approx.step3}. Let $N_{\eps, \lambda}$ be defined in \eqref{def.N}. Then, $u_2$ is quantitatively estimated as 
\begin{equation}\label{est.u2}
\begin{split}
    |u_2(x,\lambda)|
    &\le 
    \exp\big(\exp(N_{\eps,\lambda})\big) 
    \|v\|_{L^2(D)} 
    e^{(\Re\sqrt{\lambda})|x|}, 
\end{split} 
\end{equation}
for sufficiently large $C$ independent of $\eps, \lambda$. 
\end{proposition}

\begin{proposition}\label{prop.global.approx}
Let $D \subset \R^3$ be a bounded domain whose complement $\R^3\setminus D$ is connected. For $\lambda\in \Sigma_{\pi-\delta} \cap \{|\lambda| \ge 1\}$ with $\delta\in(0,\pi/2)$, let $v\in H^2(D)^3$ and $q\in H^1(D)$ satisfy
\[
    \left\{
    \begin{array}{ll}
    \lambda v - \Delta v + \nabla q
    = 0&\mbox{in}\ D, \\
    \opdiv v
    = 0&\mbox{in}\ D. 
    \end{array}\right.
\]
Then, for any $\eps>0$, there exists $(u,p)$ solving 
\[
    \left\{
    \begin{array}{ll}
    \lambda u - \Delta u + \nabla p
    = 0&\mbox{in}\ \R^3, \\
    \opdiv u
    = 0&\mbox{in}\ \R^3 
    \end{array}\right.
\]
such that $u$ approximates $v$ in $D$ as 
\[
    \|v - u\|_{L^2(D)} 
    \le 
    \eps 
    \Big( \|v\|_{H^{1}(D)} + \|q\|_{H^{1}(D)} \Big). 
\]
Moreover, $u$ can be decomposed into $u = u_1 + u_2$, where $u_1$ is a solution to the Stokes resolvent problem in $\R^3$ with polynomial growth, and $u_2$ is a solution to the heat resolvent problem in $\R^3$ with exponential growth; see Proposition \ref{prop.global.approx.step3}. Their quantitative estimates are provided in Propositions \ref{prop.global.approx.step4} and \ref{prop.global.approx.step5}, respectively.
\end{proposition}

\begin{proofx}{Proposition \ref{prop.global.approx}}
The assertion is a combination of Propositions \ref{prop.global.approx.step1}--\ref{prop.global.approx.step5}. 
\end{proofx}

\end{document}
